\documentclass[12pt]{article}
\usepackage[utf8]{inputenc}
\usepackage[T1]{fontenc}
\usepackage{hyperref}
\title{Malasar Language Resources for K-12 Curriculum Development}
\author{Compiled by Claude Sonnet 5.5 (claude-sonnet-5-5) for Andr\'as Kornai}
\date{}
\begin{document}
\maketitle

\section{Introduction}
Malasar (ISO 639-3: ymr; Glottocode mala1458) is a small Southern Dravidian language of the Malasar people. They live in Kerala and Tamil Nadu, mainly around Parambikulam and the Anamalai hills. Wikipedia gives about 7,800 speakers (2001 census). It should not be confused with Mala Malasar (ISO ima), a separate Tamiloid language. One SPPEL page says Malasar speakers use Tamil or Malayalam with outsiders.

\section{Primary Educational Resources}
No primers, textbooks or school programmes for Malasar were found in the searches.

\section{Linguistic and Reference Materials}
\begin{itemize}
\item Wikipedia, ``Malasar language'': \url{https://en.wikipedia.org/wiki/Malasar_language}
\item SPPEL (Scheme for Protection and Preservation of Endangered Languages) documentation page. It says documentation is in progress, with a grammar, a trilingual Malasar-English-Tamil electronic dictionary and an ethnolinguistic profile planned. I found no evidence that these have been published. \url{https://sppel.org/malasardoc.aspx}
\item SPPEL language documentation resources for Malasar: \url{https://sppel.org/resources/language-documentation/malasar}
\item Max Planck EVA numerals page for Malasar: \url{https://lingweb.eva.mpg.de/channumerals/Malasar.htm}
\end{itemize}

\section{Religious and Community Resources}
\begin{itemize}
\item Wikipedia, ``Malasar people'', on the community: \url{https://en.wikipedia.org/wiki/Malasar_people}
\end{itemize}

\section{Digital and Audio Resources}
\begin{itemize}
\item CIIL Sanchika special collection, which appeared in the results. I did not verify that it contains Malasar material. \url{https://sanchika.ciil.org/communities/a4c3ec04-9357-4511-8858-89098b27bceb}
\end{itemize}

\section{Curriculum Development Context}
The SPPEL project is the most relevant effort, and its outputs should be requested from SPPEL and CIIL directly. Tamil and Malayalam are the school languages in the Malasar region. Malasar-medium materials would have to be developed with the community, building on the planned SPPEL dictionary and grammar.
\end{document}